\chapter{Lecture 8: Preprocessor, Separate Compilation and Make}
\begin{center}
\textit{January 30, 2018}
\end{center}

\section{The Preprocessor (continued)}
Last time: \code{\#include}, and \code{g++ -E -P file.cc} to see the
preprocessor output. Also \code{\#define VAR VALUE}:
\begin{lstlisting}[style=cpp]
#define MAX 10
...
int arr[MAX];
\end{lstlisting}

\subsection{Conditional Compilation}
\begin{lstlisting}[style=cpp]
#define IOS 1
#define BBOS 2
#define OS IOS

#if OS == BBOS
    long long int key;     // only one of these lines
#elif OS == IOS
    int key;               // survives preprocessing
#endif
\end{lstlisting}
This requires manually editing the definition of \code{OS} to target IOS or
BlackBerry OS. Better: define preprocessor variables from the command line.
\begin{lstlisting}[style=bash]
g++ -DOS=2 file.cc      # same as putting "#define OS 2" at the top
\end{lstlisting}

A handy trick for ``commenting out'' a large block of code (which may itself
contain \code{/* */} comments):
\begin{lstlisting}[style=cpp]
#if 0
    // block of code the compiler never sees
#endif
\end{lstlisting}

\subsection{Defined or Not}
\code{\#define VAR} with no value defines \code{VAR} with the empty string as its
value.
\begin{itemize}
  \item \code{\#ifdef VAR} --- true if \code{VAR} has been defined.
  \item \code{\#ifndef VAR} --- true if \code{VAR} has \emph{not} been defined.
\end{itemize}
\begin{lstlisting}[style=cpp]
// preprocess/debug.cc
int main() {
#ifdef DEBUGGING
    cerr << "setting x = 1" << endl;
#endif
    int x = 1;
    ...
}
\end{lstlisting}
\begin{lstlisting}[style=bash]
g++ -DDEBUGGING debug.cc   # debug output on
g++ debug.cc               # debug output compiled out
\end{lstlisting}

\section{Separate Compilation}
Split a program into modules, each with:
\begin{itemize}
  \item a \textbf{header} (\code{.h}) file: type definitions and declarations
        (prototypes) of functions --- the \emph{interface};
  \item an \textbf{implementation} (\code{.cc}) file: the definitions
        (bodies) of the functions --- the \emph{implementation}.
\end{itemize}

\begin{lstlisting}[style=cpp]
// vec.h
struct Vec {
    int x, y;
};
Vec operator+(const Vec &v1, const Vec &v2);

// vec.cc
#include "vec.h"
Vec operator+(const Vec &v1, const Vec &v2) {
    return {v1.x + v2.x, v1.y + v2.y};
}

// main.cc
#include "vec.h"
int main() {
    Vec v{1, 2};
    v = v + v;
}
\end{lstlisting}

\begin{definition}
A \textbf{declaration} asserts that something exists (e.g.\ a function
prototype). A \textbf{definition} gives its full details and allocates space
(e.g.\ a function body). An entity can be declared many times but defined only
once.
\end{definition}

\subsection{How to Compile}
\begin{enumerate}
  \item \code{g++ *.cc} --- only works if every \code{.cc} file in the
        directory is part of the project.
  \item List all the files: \code{g++ vec.cc main.cc}.
  \item Compile separately. \code{g++ main.cc} on its own fails: by default
        \code{g++} tries to produce an executable, so it tries to \emph{link},
        and \code{operator+} is missing. Use \code{-c} to compile only:
\begin{lstlisting}[style=bash]
g++ -std=c++14 -c vec.cc     # produces vec.o
g++ -std=c++14 -c main.cc    # produces main.o
g++ main.o vec.o -o myprog   # link the object files
\end{lstlisting}
\end{enumerate}

\begin{important}[Never compile headers]
\code{.h} files are meant to be \emph{included}, never compiled.
\end{important}

\code{-c} produces \textbf{object files} (\code{.o}). An object file contains
the compiled binary code, plus a list of what the file \emph{provides} and what
it still \emph{needs} from other files. Linking matches these up. \code{-o name}
names the output.

The benefit: if only \code{vec.cc} changes, we only recompile \code{vec.cc}
and re-link.

\section{Make}
The problem: keeping track of which files need to be recompiled after a change
(and remembering that changing \code{vec.h} means recompiling both \code{.cc}
files). The tool for this is \textbf{make}.

Create a file called \code{Makefile} that describes the layout of the program
and the dependencies between files:
\begin{itemize}
  \item \code{myprog} depends on \code{main.o} and \code{vec.o};
  \item \code{vec.o} depends on \code{vec.cc} and \code{vec.h};
  \item \code{main.o} depends on \code{main.cc} and \code{vec.h}.
\end{itemize}

\begin{lstlisting}[style=bash]
# tools/make/example1/Makefile
# target: dependencies
#<TAB>recipe (the indentation MUST be a tab character)
myprog: main.o vec.o
	g++ main.o vec.o -o myprog

vec.o: vec.cc vec.h
	g++ -std=c++14 -c vec.cc

main.o: main.cc vec.h
	g++ -std=c++14 -c main.cc

.PHONY: clean
clean:
	rm *.o myprog
\end{lstlisting}

\begin{itemize}
  \item \code{make} builds the first target (\code{myprog}); \code{make clean}
        builds the \code{clean} target. The first time, everything is compiled.
  \item If \code{vec.cc} changes, make sees that the timestamp on \code{vec.cc}
        is newer than \code{vec.o} $\Rightarrow$ recompile \code{vec.o}. Now
        \code{vec.o} is newer than \code{myprog} $\Rightarrow$ relink
        \code{myprog}. \code{main.o} is untouched.
  \item The \code{clean} target deletes all binaries so you can start fresh.
        \code{.PHONY} tells make that \code{clean} is not the name of a file.
\end{itemize}

\subsection{Generalizing with Variables}
\begin{lstlisting}[style=bash]
# example2
CXX = g++
CXXFLAGS = -std=c++14 -Wall
OBJECTS = main.o vec.o
EXEC = myprog

${EXEC}: ${OBJECTS}
	${CXX} ${CXXFLAGS} ${OBJECTS} -o ${EXEC}

main.o: main.cc vec.h
vec.o: vec.cc vec.h

.PHONY: clean
clean:
	rm ${OBJECTS} ${EXEC}
\end{lstlisting}
For a rule of the form \code{x.o: x.cc a.h b.h}, we can leave out the recipe:
make infers \code{\$\{CXX\} \$\{CXXFLAGS\} -c x.cc -o x.o}. (\code{-Wall} turns
on all warnings --- always use it.)

\subsection{Letting g++ Track Dependencies}
Keeping the dependency lists correct by hand is error prone. \code{g++} can
figure them out:
\begin{lstlisting}[style=bash]
g++ -MMD -c vec.cc
# -c   produces vec.o
# -MMD also produces vec.d containing:  vec.o: vec.cc vec.h
\end{lstlisting}

\begin{lstlisting}[style=bash]
# example3 -- the general-purpose Makefile
CXX = g++
CXXFLAGS = -std=c++14 -Wall -MMD
OBJECTS = main.o vec.o
DEPENDS = ${OBJECTS:.o=.d}
EXEC = myprog

${EXEC}: ${OBJECTS}
	${CXX} ${CXXFLAGS} ${OBJECTS} -o ${EXEC}

-include ${DEPENDS}

.PHONY: clean
clean:
	rm ${OBJECTS} ${EXEC} ${DEPENDS}
\end{lstlisting}
\code{-include \$\{DEPENDS\}} pastes in the generated \code{.d} files (the
\code{-} means ``don't complain if they don't exist yet''). For a new project,
only \code{OBJECTS} and \code{EXEC} need to change.

\section{Include Guards}
Suppose \code{main.cc} includes \code{vec.h} and also includes some other
header (say \code{linalg.h}) which itself includes \code{vec.h}:
\begin{lstlisting}[style=bash]
$ g++ -c main.cc
error: redefinition of 'struct Vec'
\end{lstlisting}
The struct was pasted in twice. The fix is an \textbf{include guard}:
\begin{lstlisting}[style=cpp]
// vec.h
#ifndef VEC_H
#define VEC_H
struct Vec {
    int x, y;
};
Vec operator+(const Vec &v1, const Vec &v2);
#endif
\end{lstlisting}
The first time the file is included, \code{VEC\_H} is not defined, so it gets
defined and the contents are kept. Every later inclusion sees \code{VEC\_H}
already defined and skips everything.

\begin{important}[Rule]
Always put include guards in all \code{.h} files.
\end{important}

\section*{Summary}
\begin{itemize}
  \item \code{\#if}/\code{\#ifdef}/\code{\#ifndef} and \code{-DVAR=VALUE} give
        conditional compilation.
  \item Split code into headers (interface) and \code{.cc} files
        (implementation); compile each with \code{-c} and link the \code{.o}
        files.
  \item \code{make} rebuilds only what is out of date, based on timestamps; use
        variables and \code{-MMD} to get a reusable Makefile.
  \item Every header gets an include guard.
\end{itemize}
